\scanpage{392}\mspage{\deleted{\unsure{130}}\unsure{131|130}}

Cela signifie, à très peu de choses près, qu'on a une interprétation de $\underline{\mathrm{Ch}}_{\infty,\overline\Sigma}(a,b)$, en tant qu'objet de $\mathrm{Hot}$, comme un ``$\mathrm{Hom}$ externe'' \inserted{$\mathrm{Hom\,ext}(a,b)$} à valeurs dans $\mathrm{Hot}$, défini sur la catégorie \inserted{$M=$} $X$ à l'aide du seul ensemble admissible \deleted{\unsure{\Sigma}} de flèches $\overline\Sigma\subset\mathrm{Fl}(\deleted{X}M)$, \inserted{opération} satisfaisant la condition

\[
\left|\enspace\underline{\mathrm{Hom}_{\mathrm{Hot}}(Y,\deleted{\unsure{\mathrm{Ch}}}\mathrm{Hom\,ext}(a,b))\simeq\mathrm{Hom}_{\deleted{X}M(Y)\overline\Sigma_Y^{-1}}(a_X,b_X)}\enspace\right|
\tag*{(*)}
\]

(qui peut s'interpréter aussi, soit comme

\[
\mathrm{Hom}_{\deleted{X}M}(p_!(a_{\unsure{Y|X}}),b)
\]

si $p_!:\underbrace{M(Y)\overline\Sigma_Y^{-1}}_{\mathbb D(Y)}\inserted{\longrightarrow M}$ existe, soit comme

\[
\mathrm{Hom}_{M}(a_{\deleted{\unsure{X|Y}}},p_*(b_{\unsure{Y|X}}))
\]

si $p_*:\underbrace{M(Y)\overline\Sigma_Y^{-1}}_{\mathbb D Y}\longrightarrow\underset{\substack{\Vert\\ \mathbb D(e)}}{M}$ existe)\,; avec le seul petit ennui que cette formule n'est valable que pour $Y$ \uline{\unsure{fini|finie}}. Si on veut exprimer le second membre pour $Y$ variable quelconque dans $\mathrm{Cat}$, il ne suffit pas d'un objet de $\mathrm{Hot}$ ($\mathrm{Hom\,ext}(a,b)$) dans le premier membre, mais il faut le Ind-objet $H(a,b)=(\mathrm{Ch}_n(a,b))_{n\in\mathbb N}$ dans $\mathrm{Hot}$. Si donc on tient à travailler \deleted{avec des catégories} \inserted{dans un ``domaine''} \inserted{$\mathrm{Diag}$} du dérivateur $\mathbb D=\mathbb D_{M,\unsure{\Sigma|\overline\Sigma}}$ aussi grand que possible, en se bornant \deleted{\unsure{pour}} à des catégories finies

% archive sheet number: 391
